\begin{afterword}
\summaryhead

The author argues that science as a social process is slow. Scientists change their minds for bad reasons as well as good ones, and each unjustified shift must be undone with extra evidence. His examples are the gradual acceptance of many-worlds, which he attributes to physicists waiting for a large enough ``pack,'' and debates over evolutionary psychology.

He then asks whether individuals can reach correct conclusions faster than science. Science says nothing about where hypotheses come from. But in a large space of possibilities, such as all equations that fit in 32 bits, most of the evidence is spent finding the right hypothesis, not confirming it. A scientist who proposes the right equation must already have done most of the work. In the interval before the experiment, the scientist rationally knows something science has not yet confirmed, and this interval is the frontier of science.

\responsehead

The title promises individuals who outrun science. The essay delivers an account of one stage inside science, the proposing of hypotheses. It treats this stage as neglected, when philosophers from Peirce to Reichenbach to Popper have discussed it for over a century. Peirce's 1903 remark about ``trillions of trillions of hypotheses'' is the essay's central idea, uncredited.

The argument for the title turns on one word. By the essay's own arithmetic, a scientist who has raised a hypothesis to 10\% has done most of the work. The essay then says that scientist ``rationally knew'' something. A belief at 10\% is a guess that fails nine times in ten, and the essay concedes the scientist ``may not \emph{know}'' the work was done. It reasons only from the scientist who turned out right and ignores those who reasoned the same way and were wrong. Knowledge that cannot be told from a lucky guess until the experiment is run is what Science's distrust is for. The individual's guess still waits on science to become knowledge.

The bits the essay credits to the individual come mostly from science. No working scientist starts from a uniform prior over four billion equations. They start from what their training makes plausible, and that training is the product of the slow public process they are said to outrun.

The evidence for science's inefficiency is the author's own disputes. Many-worlds is said to be gaining ground, with no data. Seventeen years later, a 2025 \textsc{Nature} survey of more than 1,100 physicists found 15\% favoring it, well behind the Copenhagen interpretation. Collapse is said to violate relativity, when a relativistic collapse model had been published two years earlier. The evolutionary psychology debate is reduced to one side's ``ridiculous objections,'' with no name, paper or objection given. Planck's remark is quoted as a finding, though studies that tested it found mixed results. The essay never considers that slow, divided conversion can serve a community, as Kitcher argued in 1990.

The misreadings of Bayes do the most damage, in an essay that asks readers to prefer Bayes to science. Bayes's rule is not ``falsified'' by outcomes, and it would not ``burst into flames'' at a reliable Ouija board. It would update on it.

In short: the old observation that someone must guess before anyone tests, renamed knowledge, and used to certify the author's positions as ahead of science.
\end{afterword}
